\clearpage
\begingroup
\fancyhead[L]{\small NACHLASS UND BRIEFWECHSEL}
\renewcommand{\headrulewidth}{0pt}
\par\medskip
\Needspace{8\baselineskip}
\pdfbookmark[2]{[3]}{pallas-IX-fragment3}
\begin{center}[3.]\end{center}
\begingroup
\small
\newcommand{\IXmars}{\text{\AstroSymbols ♂}}
\newcommand{\IXpallas}{\text{\AstroSymbols ⚴}}
\par\noindent\hspace*{\parindent}[Durch Multiplikation mit den Faktoren der Gleichungen 2) erhält man schließlich die Ausdrücke für
\(\frac{\mathrm{d}\log\mathrm{hyp}\,n}{n\,\mathrm{d}t}\) und \(\frac{\mathrm{d}\varphi}{n\,\mathrm{d}t}\). Der erste Ausdruck liefert bei einer Integration wie auf den Seiten 523---524
die Störungen des Logarithmus der großen Halbachse und des ersten Teils der Epoche. Der zweite Ausdruck liefert
die Störungen von \(\varphi\). Man erhielt Folgendes.
\[
\begin{alignedat}{2}
\delta\log a&=1\sin(\IXmars-3\IXpallas+10^0\ 27'\ 49'')
&\mspace{4mu}\delta\!\int n\,\mathrm{d}t&=0\overset{\prime\prime}{,}18\sin(\IXmars-3\IXpallas+280^0\ 27'\ 49'')\\
&\quad+1\sin(\IXmars-2\IXpallas+202\ 20\ 18)
&&\enspace+0{,}16\sin(\IXmars-2\IXpallas+292\ 20\ 18)\\
&\quad+6\sin(\IXmars-\IXpallas+90\ 54\ 56)
&&\enspace+0{,}28\sin(\IXmars-\IXpallas+180\ 54\ 56)\\
&\quad+11\sin(2\IXmars-5\IXpallas+124\ 17\ 38)
&&\enspace+8{,}14\sin(2\IXmars-5\IXpallas+34\ 17\ 38)\\
&\quad+2\sin(2\IXmars-4\IXpallas+162\ 24\ 56)
&&\enspace+0{,}12\sin(2\IXmars-4\IXpallas+252\ 24\ 56)\\
&\quad+2\sin(3\IXmars-7\IXpallas+76\ 26\ 28)
&&\enspace+0{,}38\sin(3\IXmars-7\IXpallas+166\ 26\ 28)\\
&\quad+2\sin(4\IXmars-10\IXpallas+171\ 18\ 55)
&&\enspace+0{,}79\sin(4\IXmars-10\IXpallas+81\ 18\ 55)\\
&\quad+1\sin(5\IXmars-12\IXpallas+120\ 44\ 45)
&&\enspace+0{,}28\sin(5\IXmars-12\IXpallas+210\ 44\ 45).
\end{alignedat}
\]
\[
\begin{aligned}
\delta\varphi&=0\overset{\prime\prime}{,}15\sin(\IXmars+335^0\ 22'\ 13'')\\
&\quad+0{,}49\sin(2\IXmars-5\IXpallas+106\ 33\ 27).
\end{aligned}
\]
\begin{center}
\begin{tabular}{@{}lrl@{}}
Säkularer Teil von \(\delta\varphi\): & täglich & \(-0\overset{\prime\prime}{,}000\,0059\)\\
& jährlich & \(-0\overset{\prime\prime}{,}00215\).\\
\end{tabular}
\end{center}
\par\noindent\hspace*{\parindent}\(\delta\log a\) ist in Einheiten der 7. Dezimalstelle angegeben. Berücksichtigt sind alle Gleichungen, deren Beträge
eine solche Einheit bzw.\ \(0\overset{\prime\prime}{,}10\) erreichen.]
\par\smallskip\noindent VII.\hfill 75\par
\par\endgroup

\par\medskip
\SourceLeaf{600}
\Needspace{8\baselineskip}
\pdfbookmark[2]{[4]}{pallas-IX-fragment4}
\begin{center}[4.]\end{center}
\begingroup
\small
\par\noindent\hspace*{\parindent}[Die Reihen, in die oben die Größen \(\frac{x'-x}{\rho^3}\), \(\frac{y'-y}{\rho^3}\), \(\frac{z'-z}{\rho^3}\) entwickelt wurden,
konvergieren so langsam, dass man trotz der kleinen Beträge der einzelnen Störungsgleichungen 24 Positionen des Mars
und ebenso viele Positionen der Pallas der Rechnung zugrunde legen musste. Man kann hier auch anders vorgehen.
Dazu multipliziert man die zu entwickelnden Größen mit einem Faktor \(\lambda\),
der so gewählt wird, dass die Entwicklungen der Produkte \(\frac{x'-x}{\rho^3}\lambda\) usw. möglichst gut konvergieren.
Die langsam konvergierende Entwicklung von \(\lambda^{-1}\) lässt sich mithilfe analytischer
Formeln leicht beliebig weit fortsetzen. Die Produkte \(\frac{x'-x}{\rho^3}\lambda\) usw. kann man dann leicht durch Interpolation entwickeln,
wobei nur wenige Marspositionen zugrunde gelegt werden. Die Entwicklung von \(\lambda^{-1}\) lässt sich mit
wenig Aufwand so weit fortführen, wie es nötig erscheint. Multipliziert man beide Entwicklungen, so erhält man
die gesuchte Entwicklung von \(\frac{x'-x}{\rho^3}\) usw. Diese Methode kann auch dazu dienen, die allein durch Interpolation
gewonnenen Entwicklungen zu überprüfen.
\par\noindent\hspace*{\parindent}Der Faktor \(\lambda\) wurde folgendermaßen bestimmt.]
\par\endgroup
\clearpage
\endgroup
